\documentclass{article}
\usepackage[margin=1.5cm]{geometry}
\usepackage{amsmath}
\usepackage{enumitem}

\setlist{nosep,leftmargin=*}

\begin{document}
\twocolumn

\title{Chapter 7 In-Class Exercise: Fourier Series as a Module}
\author{Prof. Jordan C. Hanson}
\date{}

\maketitle

\section{Step 1: Create the Module}

\begin{enumerate}

\item Previously, we created separate functions for one Fourier term, the Fourier sum, and the plot.  We will now place those functions inside a module. Create a file called

\begin{verbatim}
fourier_tools.py
\end{verbatim}

and begin with

\begin{verbatim}
import numpy as np
import matplotlib.pyplot as plt
\end{verbatim}

Add the function that evaluates the $n$-th term:

\[
a_n(x)=
2\frac{(-1)^{n+1}}{n}\sin(nx).
\]

\begin{verbatim}
def fourier_term(x, n):
    term = 2 * (-1)**(n + 1)
    term = term * np.sin(n*x) / n
    return term
\end{verbatim}

\end{enumerate}


\section{Step 2: Add the Sum}

\begin{enumerate}

\item Add a second function to the same module:

\begin{verbatim}
def sawtooth(x, N):
    total = 0.0
    for n in range(1, N + 1):
        total = total + \
                fourier_term(x, n)
    return total
\end{verbatim}

Notice that \texttt{sawtooth()} calls another function located inside the same module.

\begin{itemize}

\item Which function performs one Fourier-series term?

\item Which function controls how many terms are added?

\end{itemize}

\end{enumerate}


\section{Step 3: Add the Plot Function}

\begin{enumerate}

\item Add a third function:

\begin{verbatim}
def plot_series(x, N):
    y = sawtooth(x, N)
    plt.plot(x, y, label=f"N = {N}")
\end{verbatim}

The module now contains three related functions:

\begin{verbatim}
fourier_term()
sawtooth()
plot_series()
\end{verbatim}

Save the completed file.

\end{enumerate}


\section{Step 4: Upload and Import}

\begin{enumerate}

\item In Google Colab, upload \verb+fourier_tools.py+. In a new code cell, import the module:

\begin{verbatim}
import fourier_tools
\end{verbatim}

Also import the libraries needed by the main program:

\begin{verbatim}
import numpy as np
import matplotlib.pyplot as plt
\end{verbatim}

Test one function directly:

\begin{verbatim}
print(fourier_tools.fourier_term(np.pi/4, 1))
\end{verbatim}

\end{enumerate}


\section{Step 5: Main Program}

\begin{enumerate}

\item Write the following script in a Colab code cell:

\begin{verbatim}
x = np.linspace(-np.pi,np.pi, 400)
fourier_tools.plot_series(x, 1)
fourier_tools.plot_series(x, 5)
fourier_tools.plot_series(x, 20)
plt.plot(x, x, label="f(x) = x")
plt.xlabel("x")
plt.ylabel("f(x)")
plt.title("Sawtooth Fourier Series")
plt.legend()
plt.grid()
plt.show()
\end{verbatim}

Run the program and inspect the graph.

\begin{itemize}

\item Which file contains the mathematical functions?

\item Which part of the program determines $N$?

\item Change the final value from $N=20$ to $N=50$ and run the program again.

\end{itemize}

\end{enumerate}


\section{Challenge}

Modify the main program so that the values of $N$ are stored in \verb+terms = [1, 5, 20, 50]+ and use a \texttt{for}  loop to call \texttt{fourier\_tools.plot\_series()} for each value.

\end{document}